\chapter{European Equity Market Structure}\label{ch:m1:european-equity-market-structure}

A share of a large French company trades this morning on the Paris exchange
that listed it, on several pan-European platforms, in a handful of banks'
own books, in auctions that last a tenth of
a second, and in \omterm{def:m1:us-equity-market-structure:dark}{dark pools} that match at a midpoint nobody publishes in one
place. In July 2026 about \euro 80 billion of European shares changed hands
each day; barely half of what traded on venues did so in a continuous,
visible order book. And until September 2026 the continent had never had what
the United States has long had: a single public tape saying what the best
price is. This chapter describes a market that solved the same
problems as the American one with opposite choices.

\section{One directive, four kinds of place}

The Markets in Financial Instruments Directive ended the national exchanges'
monopolies in 2007; its second version, MiFID II, with the accompanying
regulation MiFIR, has governed European trading since January 2018. It sorts
every place where a share can trade into legal categories.

\begin{definition}[Regulated market and organised trading facility]\label{def:m1:european-equity-market-structure:rm}
A \emph{regulated market}\index{regulated market} is a multilateral system,
operated by a market operator, that brings together multiple third-party
buying and selling interests under non-discretionary rules and admits
instruments to trading: the European legal form of an exchange. An
\emph{organised trading facility}\index{organised trading facility} (OTF) is
a multilateral system whose operator may exercise discretion in matching; it
exists for bonds and derivatives and may not trade shares.
\end{definition}

Between the two sits the \omterm{def:m1:exchanges-brokers-venues:mtf}{multilateral trading facility}
(\cref{def:m1:exchanges-brokers-venues:mtf}), which has the \omterm{def:m1:european-equity-market-structure:rm}{regulated
market}'s rules without its listing function: the pan-European platforms are
MTFs.

\begin{definition}[Systematic internaliser]\label{def:m1:european-equity-market-structure:si}
A \emph{systematic internaliser}\index{systematic internaliser} (SI) is an
investment firm that, on an organised, frequent and substantial basis, deals
on its own account when executing client orders outside a trading venue. It
is a counterparty, not a venue: every trade is bilateral, against the firm's
own book, and the firm must publish firm quotes in liquid shares up to a
standard size.
\end{definition}

\begin{omfigure}
\begin{tikzpicture}[font=\small,
  top/.style={draw=omExo,rounded corners=2pt,minimum width=3.4cm,minimum height=0.8cm,align=center,fill=omExo!6},
  box/.style={draw=omDef,rounded corners=2pt,minimum width=2.9cm,minimum height=1.25cm,align=center,fill=omDef!4},
  bil/.style={draw=omThm,rounded corners=2pt,minimum width=2.9cm,minimum height=1.25cm,align=center,fill=omThm!5}]
\node[top] (m) at (4.0,2.2) {Multilateral\\(many-to-many)};
\node[top] (b) at (12.0,2.2) {Bilateral\\(firm against client)};
\node[box] (rm) at (0.6,0) {Regulated market\\lists and trades};
\node[box] (mtf) at (4.0,0) {MTF\\trades only};
\node[box] (otf) at (7.4,0) {OTF\\not for shares};
\node[bil] (si) at (10.4,0) {Systematic\\internaliser};
\node[bil] (otc) at (13.6,0) {Other OTC};
\foreach \x in {rm,mtf,otf}{ \draw[omExo,thick] (m.south) -- (\x.north); }
\foreach \x in {si,otc}{ \draw[omExo,thick] (b.south) -- (\x.north); }
\end{tikzpicture}

\omcaption{Where a European share may legally trade. A \emph{share trading
obligation} requires investment firms to execute shares on a \omterm{def:m1:european-equity-market-structure:rm}{regulated
market}, on an MTF or with a \omterm{def:m1:european-equity-market-structure:si}{systematic internaliser}, which confines true
over-the-counter trading to exceptional cases.}\label{fig:m1:european-equity-market-structure:taxonomy}
\end{omfigure}

\begin{definition}[Best execution]\label{def:m1:european-equity-market-structure:bestex}
\emph{Best execution}\index{best execution} is the obligation of a firm
executing client orders to take all sufficient steps to obtain the best
possible result, taking into account price, costs, speed, likelihood of
execution and settlement, size and nature of the order.
\end{definition}

This is the European answer to fragmentation, and it is the opposite of the
American one. There is no \omterm{def:m1:us-equity-market-structure:opr}{order protection rule}
(\cref{def:m1:us-equity-market-structure:opr}): no venue is obliged to route
to a better price elsewhere, and a trade at an inferior price breaks no
market rule. The duty sits with the \emph{broker}, as a policy it must write,
follow and be able to demonstrate. Venues connect to whom they like; brokers'
routers are what ties the market together.

\section{Fragmentation without a tape}

\begin{definition}[Consolidated tape provider]\label{def:m1:european-equity-market-structure:ctp}
A \emph{consolidated tape provider}\index{consolidated tape provider} (CTP)
is an entity authorised to collect trade reports, and for shares the best
bids and offers, from all venues and publication arrangements, and to
disseminate them as one continuous stream.
\end{definition}

\begin{dated}{2026-09}{A tape, at last}\label{dat:m1:european-equity-market-structure:ctp}
For eighteen years after competition began, Europe had no consolidated tape
for shares: each firm built its own view from venue feeds, or bought a
vendor's. In December 2025 ESMA selected a consortium of exchanges, EuroCTP,
as the first provider for shares and exchange-traded funds, and later
authorised it for five years under its direct supervision. The tape's launch
was set for 14 September 2026, with a transition period to 30 September
2026. Its data are free of charge for retail investors, academics and
regulators.
\end{dated}

The absence of an official best price had practical consequences that
persist. ``The European best bid and offer'' is a number each firm computes
for itself, from the venues it has chosen to connect to, with no round-lot
filter and no protected status; two brokers can honestly disagree about it.
Transaction-cost analysis must state which venues its benchmark included.
And the primary exchange's price keeps a privileged role out of habit: when
the national exchange halts, volume on the platforms that copy its price
often dries up too.

\begin{definition}[Effective number of venues]\label{def:m1:european-equity-market-structure:hhi}
If venues have market shares $s_1,\dots,s_n$ summing to one, the
\emph{Herfindahl index} is $H = \sum_i s_i^2$ and the \emph{effective number
of venues} is $1/H$: the number of equal-sized venues that would produce the
same concentration.
\end{definition}

\begin{example}[Fragmented, but how much?]\label{ex:m1:european-equity-market-structure:hhi}
A stock trades 60\% on its primary exchange, 20\% on one platform and 10\% on
each of two others. $H = 0.36 + 0.04 + 0.01 + 0.01 = 0.42$: an effective 2.4
venues, although there are four. If the primary's share falls to 40\% and the
others take 20\% each, $H = 0.28$ and the effective number is 3.6. The measure
moves mostly with the leader's share
(\cref{fig:m1:european-equity-market-structure:effective}).
\end{example}

\begin{omfigure}
\begin{tikzpicture}
\begin{axis}[width=10.5cm,height=4.6cm,
  xlabel={market share of the largest venue (\%)},ylabel={effective number of venues},
  tick label style={font=\small},label style={font=\small},
  xmin=20,xmax=100,ymin=1,ymax=6,grid=major,grid style={gray!25},table/col sep=comma]
\addplot[omDef,very thick] table[x=leader_pct,y=effective]{figdata/markets-1/11-european-equity-market-structure/effective.csv};
\end{axis}
\end{tikzpicture}

\omcaption{Effective number of venues when one venue has the share on the
horizontal axis and five others split the rest equally. Six venues are
``six'' only when the leader is down to a sixth of the market. Data: computed
by the chapter's script.}\label{fig:m1:european-equity-market-structure:effective}
\end{omfigure}

\section{The mechanisms: lit, dark, periodic, close}

\begin{dated}{2026-09}{How European shares traded in July 2026}\label{dat:m1:european-equity-market-structure:july}
By one pan-European operator's monthly statistics, on-exchange turnover in
European equities averaged \euro 57.5 billion a day in July 2026, out of
\euro 80 billion of addressable activity including off-exchange trading. Of
the on-exchange part, continuous central limit order books were 52.8\%,
closing auctions 24.3\%, \omterm{def:m1:european-equity-market-structure:periodic}{periodic auctions} 10.4\% and non-displayed order
books 10.9\%. \omterm{def:m1:european-equity-market-structure:si}{Systematic internalisers} accounted for \euro 14.3 billion a
day, 17.8\% of addressable volume.
\end{dated}

\begin{omfigure}
\begin{tikzpicture}
\begin{axis}[width=9cm,height=4.2cm,xbar,bar width=8pt,
  xlabel={share of on-exchange turnover, July 2026 (\%)},
  ytick=data,yticklabels from table={figdata/markets-1/11-european-equity-market-structure/mechanisms.csv}{label},
  yticklabel style={font=\small},tick label style={font=\small},label style={font=\small},
  y dir=reverse,xmin=0,xmax=65,enlarge y limits=0.2,grid=major,grid style={gray!25},
  nodes near coords,nodes near coords style={font=\footnotesize},
  table/col sep=comma]
\addplot[fill=omDef!60,draw=omDef] table[x=pct,y=k]{figdata/markets-1/11-european-equity-market-structure/mechanisms.csv};
\end{axis}
\end{tikzpicture}

\omcaption{Four ways of matching on European venues. A quarter of
on-exchange turnover happens in the few minutes of the closing auctions.
Data: as in \cref{dat:m1:european-equity-market-structure:july}.}
\end{omfigure}

MiFID II requires pre-trade transparency: a venue must publish its bids and
offers. Dark trading exists by \emph{waiver} from that rule.

\begin{definition}[Reference price waiver and large-in-scale waiver]\label{def:m1:european-equity-market-structure:waivers}
Under the \emph{reference price waiver}\index{reference price waiver} a venue
need not display orders that execute at a price taken from another market,
in practice the midpoint of the primary exchange's quote. Under the
\emph{large-in-scale waiver}\index{large-in-scale waiver} it need not display
an order that is large compared with normal market size for that share.
\end{definition}

\begin{definition}[Volume cap]\label{def:m1:european-equity-market-structure:cap}
The original \emph{double volume cap}\index{double volume cap} suspended
trading of a share under the \omterm{def:m1:european-equity-market-structure:waivers}{reference price waiver} when, over twelve months,
such trading exceeded 4\% of the share's total volume on any one venue or 8\%
across the Union. Since October 2025 it is replaced by a \emph{single volume
cap}: 7\% across the Union, with a three-month suspension of the waiver in
that share when it is exceeded.
\end{definition}

\begin{dated}{2026-09}{The cap in operation}\label{dat:m1:european-equity-market-structure:svc}
ESMA published the first single-volume-cap results on 9 October 2025 and
publishes calculations and suspensions quarterly; the double-volume-cap
system was scheduled for decommissioning in January 2026. Large-in-scale
trading is outside the cap.
\end{dated}

\begin{definition}[Periodic auction]\label{def:m1:european-equity-market-structure:periodic}
A \emph{periodic auction}\index{periodic auction} is a short call auction,
typically lasting a fraction of a second, that starts when two orders in the
book could match, publishes an indicative price and size during its brief
life, and uncrosses at a single price, often constrained to lie within the
best bid and offer of the reference market.
\end{definition}

\omterm{def:m1:european-equity-market-structure:periodic}{Periodic auctions} grew when the first caps suspended dark trading in many
shares: a share suspended from midpoint dark trading could still be traded at, or
near, the midpoint in an auction that is ``lit'' by the letter of the law and
gives almost as little information as a \omterm{def:m1:us-equity-market-structure:dark}{dark pool}. Every rule in this market
has produced the mechanism that routes around it; the reader who wants to
predict the next one should ask what the new single cap makes expensive.

\begin{omfigure}
\begin{tikzpicture}
\begin{axis}[width=11cm,height=5cm,
  xlabel={month},ylabel={share of the stock's volume (\%)},
  tick label style={font=\small},label style={font=\small},
  xmin=12,xmax=36,ymin=0,ymax=14,grid=major,grid style={gray!25},unbounded coords=jump,
  legend columns=3,legend style={font=\footnotesize,at={(0.5,-0.32)},anchor=north,draw=none,fill=none,/tikz/every even column/.append style={column sep=6pt}},
  table/col sep=comma]
\addplot[omThm,very thick] table[x=month,y=usage_pct]{figdata/markets-1/11-european-equity-market-structure/cap.csv};
\addplot[omExo,thick,const plot mark mid] table[x=month,y=dark_pct]{figdata/markets-1/11-european-equity-market-structure/cap.csv};
\addplot[omDef,thick,const plot mark mid] table[x=month,y=periodic_pct]{figdata/markets-1/11-european-equity-market-structure/cap.csv};
\addplot[black,dashed] coordinates {(12,7) (36,7)};
\legend{twelve-month waiver usage,monthly dark share,monthly periodic-auction share}
\end{axis}
\end{tikzpicture}

\omcaption{A simulated share whose midpoint dark trading drifts upward. When
the twelve-month usage crosses 7\% (dashed) the waiver is suspended for three
months: dark volume falls to zero and half of it reappears in \omterm{def:m1:european-equity-market-structure:periodic}{periodic
auctions}. Data: the chapter's simulation; the migration split is
illustrative.}\label{fig:m1:european-equity-market-structure:cap}
\end{omfigure}

\section{The close, and the tick}

A quarter of on-exchange turnover in the closing auction is not an accident
of taste. Index funds must trade at the closing price they are measured
against (\cref{ch:m1:the-buy-side}); the primary exchange owns that price;
and as more volume goes there, the intraday book thins, which sends more
volume there. The auction mechanism itself is the subject of
\cref{ch:m1:auctions}.

Unlike the American one-cent increment, European tick sizes are set by a
table. Under the tick-size regime (RTS~11) the minimum increment of a share
depends on two things: its price, and a \emph{liquidity band} given by its
average daily number of transactions on its most relevant market,
recalculated once a year. More liquid shares get smaller ticks at the same
price. The design aims at a spread of a few ticks everywhere; it also means
that one share's tick can change on the first Monday of April with no change
in its price, shifting queue lengths and \omterm{def:m1:what-a-trading-firm-does:market-maker}{market makers}' economics overnight.

\section{Tutorial: measuring fragmentation}\label{tut:m1:european-equity-market-structure:frag}

\begin{tutorial}
\textbf{Goal.} Build a European best bid and offer, measure concentration,
and monitor a volume cap. \textbf{End state:}
\cref{fig:m1:european-equity-market-structure:effective,fig:m1:european-equity-market-structure:cap}.
\begin{tutsteps}
\item \textbf{Concentration and the best quote.} There is no round-lot filter
and no protected status: the best price is simply the best price among the
venues you chose to look at.
\omcode{code/markets-1/11-european-equity-market-structure/python/euro_frag.py}{11}{27}{Herfindahl index, effective number of venues, and a best bid and offer across venues.}\label{lst:m1:european-equity-market-structure:hhi}
\item \textbf{Try it.} With Paris at $45.20 \times 45.24$ and two platforms
at $45.21 \times 45.23$ and $45.21 \times 45.24$, the best quote is $45.21
\times 45.23$: tighter than the primary exchange's own, and invisible to
anyone who watches only the primary.
\item \textbf{The cap is a rolling ratio.}
\omcode{code/markets-1/11-european-equity-market-structure/python/euro_frag.py}{30}{34}{Twelve-month usage of the reference price waiver.}\label{lst:m1:european-equity-market-structure:cap}
\item \textbf{Simulate a breach.} In the chapter's scenario the usage crosses
7\% in month 28; the dark book is empty in months 29 to 31 and the
periodic-auction share jumps from 8\% to about 12\%.
\end{tutsteps}
\textbf{What to change next.} Make usage a forecast: given the last eleven
months, what dark share next month triggers a suspension? A broker's router
needs exactly that number. Then recompute
\cref{ex:m1:european-equity-market-structure:hhi} counting \omterm{def:m1:european-equity-market-structure:si}{systematic
internalisers} as one more ``venue''.
\end{tutorial}

\section{Build: the fragmentation monitor}\label{bld:m1:european-equity-market-structure:monitor}

\begin{build}
\bfield{Purpose} The miniature firm's router and its cost analysis both need
to know, per instrument, where volume trades and by which mechanism, and how
close each share is to losing its dark liquidity.
\bfield{Interface} \texttt{Monitor.add(trade)} with a trade carrying
instrument, venue code (\cref{bld:m1:exchanges-brokers-venues:registry}),
mechanism (\texttt{lit}, \texttt{dark\_rpw}, \texttt{dark\_lis},
\texttt{periodic}, \texttt{auction}, \texttt{si}), value and month;
\texttt{shares(instrument, by)} with \texttt{by} in \texttt{venue} or
\texttt{mechanism}; \texttt{effective\_venues(instrument)};
\texttt{cap\_usage(instrument, month)}; \texttt{headroom(instrument, month)}
returning the waiver volume that would bring next month's usage to the cap.
\bfield{Rules} Values are integers in ledger units. Only \texttt{dark\_rpw}
counts toward the cap; large-in-scale does not. Usage needs twelve full
months: before that the monitor returns no figure, not a partial one.
\bfield{Acceptance tests} \texttt{code/firm/fragmentation/tests/}: shares sum
to one; the example's Herfindahl index; cap usage on a constant series; the
large-in-scale exclusion; headroom brings usage exactly to 7\%.
\bfield{Stretch} A suspension calendar: given quarterly publication dates,
the months in which each instrument's waiver is unavailable.
\end{build}

\begin{omsources}
\item Directive 2014/65/EU (MiFID II), article 4(1)(20)--(24) and article 27;
Regulation (EU) No 600/2014 (MiFIR) as amended in 2024.
\item ESMA, \emph{ESMA selects EuroCTP to become the first Consolidated Tape
Provider for shares and ETFs}, 19 December 2025; \emph{ESMA authorises
EuroCTP as the Consolidated Tape Provider for shares and exchange-traded
funds}, 2026.
\item ESMA, \emph{ESMA prepares for switch toward single volume cap in October
2025}; \emph{Double Volume Cap Mechanism} (discontinued).
\item Commission Delegated Regulation (EU) 2017/588 (RTS 11), tick size
regime for shares, depositary receipts and exchange-traded funds.
\item Cboe Global Markets, \emph{Market Metrics That Matter: European Equities
July Volume Briefing}, 17 August 2026.
\item C.-A.~Lehalle and S.~Laruelle (eds.), \emph{Market Microstructure in
Practice}, 2nd ed., World Scientific, 2018, chapter 1.
\end{omsources}

\section{Exercises}

\begin{exercise}[$\star$]\label{exo:m1:european-equity-market-structure:1}
Classify each as \omterm{def:m1:european-equity-market-structure:rm}{regulated market}, MTF, OTF, \omterm{def:m1:european-equity-market-structure:si}{systematic internaliser} or
other: (a) the exchange on which a German company is listed; (b) a
pan-European platform trading 3\,000 shares it did not list; (c) a bank
filling its clients' share orders from its own book all day; (d) a broker's
system matching bond orders with discretion.
\end{exercise}

\begin{exercise}[$\star$]\label{exo:m1:european-equity-market-structure:2}
A share's volume splits 45\%, 25\%, 15\%, 10\%, 5\% across five venues.
Compute the Herfindahl index and the effective number of venues.
\end{exercise}

\begin{exercise}[$\star$]\label{exo:m1:european-equity-market-structure:3}
From \cref{dat:m1:european-equity-market-structure:july}, compute in euros per
day the turnover of closing auctions, of \omterm{def:m1:european-equity-market-structure:periodic}{periodic auctions} and of
non-displayed books, and the off-exchange activity that is not \omterm{def:m1:european-equity-market-structure:si}{systematic
internalisers}'.
\end{exercise}

\begin{exercise}[$\star\star$]\label{exo:m1:european-equity-market-structure:4}
Three venues quote a share: X $18.340 \times 18.360$, Y $18.345 \times
18.355$, Z $18.350 \times 18.365$. Give the best bid and offer and their
venues. A broker connected only to X and Y buys at 18.355. Has it broken a
market rule? Has it failed its client?
\end{exercise}

\begin{exercise}[$\star\star$]\label{exo:m1:european-equity-market-structure:5}
Over the last eleven months a share traded \euro 4\,400 million in total, of
which \euro 300 million under the \omterm{def:m1:european-equity-market-structure:waivers}{reference price waiver}. Next month's total
is expected to be \euro 400 million. What waiver volume next month brings the
twelve-month usage exactly to 7\%? Express it as a share of that month's
volume.
\end{exercise}

\begin{exercise}[$\star\star$]\label{exo:m1:european-equity-market-structure:6}
A share at \euro 12.40 has a tick of \euro 0.005 and a typical spread of two
ticks. After the annual recalculation its liquidity band changes and the tick
becomes \euro 0.01. Give the spread in basis points before, and the minimum
possible spread after. What happens to the queue at the best price?
\end{exercise}

\begin{exercise}[$\star\star\star$]\label{exo:m1:european-equity-market-structure:7}
\emph{Coding.} With \texttt{simulate\_cap}, find the first month of breach for
seeds 0 to 49 and report the median. Then set the upward drift of the dark
share to zero and report how many of the fifty paths still breach.
\end{exercise}

\begin{exercise}[$\star\star\star$]\label{exo:m1:european-equity-market-structure:8}
\emph{Find the flaw.} A transaction-cost report states: ``Our algorithm
achieved an average fill 0.4 basis point better than the European best bid
and offer.'' The benchmark is built from the three venues the algorithm
trades on; the stock also trades on four other venues and with seven
\omterm{def:m1:european-equity-market-structure:si}{systematic internalisers}. Explain what the figure does and does not show, and
how to repair it.
\end{exercise}

\section{Problem: Where Did the Volume Go?}

\begin{problem}[{Weekend problem --- a share loses its dark pool}]\label{pb:m1:european-equity-market-structure:1}
A large-capitalisation share trades \euro 500 million a day on average, 21
days a month. Its turnover splits as follows: primary exchange's continuous
book 38\%, primary's closing auction 20\%, three MTFs' lit books 9\% each (27\%),
midpoint dark books under the \omterm{def:m1:european-equity-market-structure:waivers}{reference price waiver} 7.5\%, \omterm{def:m1:european-equity-market-structure:periodic}{periodic auctions}
4\%, large-in-scale blocks 3.5\%. (\omterm{def:m1:european-equity-market-structure:si}{Systematic internalisers} are ignored
here.)

\medskip\noindent\textbf{Part I --- Before.}
\begin{enumerate}
\item Give the monthly turnover and the monthly value traded in midpoint dark
books.
\item By \emph{venue operator} the shares are: primary 60\% (its book, its
auction and 2\% of the dark and periodic volume), and the rest split equally
among four platform operators. Compute the effective number of operators.
\item By \emph{mechanism}, compute the Herfindahl index over the six
categories above and the effective number of mechanisms.
\item The twelve-month waiver usage stands at 7.5\%. What happens, and for
how long?
\end{enumerate}

\medskip\noindent\textbf{Part II --- The suspension.}
During the suspension the 7.5\% that traded in midpoint dark books goes: 45\%
to \omterm{def:m1:european-equity-market-structure:periodic}{periodic auctions}, 25\% to large-in-scale blocks, 20\% to lit continuous
books (split between primary and MTFs in proportion to their lit shares), and
10\% is not traded at all.
\begin{enumerate}[resume]
\item Give the new share of each mechanism as a percentage of the
\emph{original} turnover.
\item Renormalise to the new, lower turnover and give the new percentages.
\item By how many percentage points did \omterm{def:m1:european-equity-market-structure:periodic}{periodic auctions} gain?
\item Recompute the effective number of mechanisms.
\item An institutional broker's midpoint fills saved it half the spread, 2
basis points, on the volume it traded dark, which was 12\% of its \euro 40
million a day in this share. If only the periodic-auction and block fractions
of that flow still execute at the midpoint, what does the suspension cost it
per month?
\end{enumerate}

\medskip\noindent\textbf{Part III --- Coming back.}
\begin{enumerate}[resume]
\item After three months the waiver is restored. With zero dark volume for
three months and 7.5\% a month before that, what is the twelve-month usage on
the day of restoration?
\item If dark trading resumes at 7.5\% a month, how many months until usage
exceeds 7\% again?
\item What monthly dark share keeps usage at exactly 7\% forever?
\item Why might the operators of dark books collectively prefer to stay under
that level, and why can they not easily coordinate to do so?
\end{enumerate}

\medskip\noindent\textbf{Part IV --- Judgement.}
\begin{enumerate}[resume]
\item Who is protected by the cap, in the legislator's reasoning?
\item Who bears its cost, according to question 9?
\item \omterm{def:m1:european-equity-market-structure:periodic}{Periodic auctions} display an indicative price for a fraction of a
second. In what sense are they lit, and in what sense dark?
\item Why does the \omterm{def:m1:european-equity-market-structure:waivers}{large-in-scale waiver} escape the cap?
\item With a consolidated tape now publishing the best bid and offer, which
statistics of this problem become easier to compute, and which do not
change?
\item State the \emph{named result}: the \omterm{def:m1:european-equity-market-structure:periodic}{periodic auctions}' market share
during the suspension, as a percentage of the renormalised turnover.
\item In one sentence: what is the general law of market structure this
problem illustrates?
\end{enumerate}
\end{problem}

\section{Interview questions}

\begin{interviewq}[$\star$ \iqroles{trader, researcher, developer}]\label{iq:m1:european-equity-market-structure:1}
Name the main differences between US and European equity market structure.
\end{interviewq}

\begin{interviewq}[$\star$ \iqroles{trader, bank}]\label{iq:m1:european-equity-market-structure:2}
What is a \omterm{def:m1:european-equity-market-structure:si}{systematic internaliser}, and how does it differ from a \omterm{def:m1:exchanges-brokers-venues:mtf}{multilateral
trading facility}?
\end{interviewq}

\begin{interviewq}[$\star\star$ \iqroles{researcher, trader}]\label{iq:m1:european-equity-market-structure:3}
Why is a quarter of European on-exchange volume traded in the closing
auction, and what does that do to intraday liquidity?
\end{interviewq}

\begin{interviewq}[$\star\star$ \iqroles{developer, researcher}]\label{iq:m1:european-equity-market-structure:4}
You must compute a European best bid and offer for transaction-cost analysis.
What decisions do you have to make that a US analyst does not?
\end{interviewq}

\begin{interviewq}[$\star\star$ \iqroles{trader, researcher}]\label{iq:m1:european-equity-market-structure:5}
A share is about to be suspended from midpoint dark trading under the volume
cap. How does that change the way you execute a large order in it?
\end{interviewq}

\begin{interviewq}[$\star\star\star$ \iqroles{researcher, trader}]\label{iq:m1:european-equity-market-structure:6}
A share's tick size doubles after the annual recalculation. Predict the
effect on spread, depth, queue times and a \omterm{def:m1:what-a-trading-firm-does:market-maker}{market maker}'s profitability, and
say how you would test your predictions.
\end{interviewq}
